\begin{figure}[tbp]
\centering
\begin{tikzpicture}[x=0.94cm,y=0.94cm,
 visible/.style={draw=black!75,line width=1.05pt},
 hidden/.style={draw=blue!65!black,line width=1.2pt},
 every node/.style={font=\small},
 exposure point/.style={circle,fill=red!65!black,inner sep=2pt}]
\draw[-{Latex[length=2mm]},black!55] (-2.6,3.1)--(-2.6,-.05);
\node[rotate=90,anchor=south,font=\footnotesize] at (-2.72,1.5)
 {backward time};
\begin{scope}
 \node[font=\small\bfseries] at (.15,3.65) {Genuine scattering};
 \fill[blue!4] (-1.65,-.15) rectangle (1.6,1.82);
 \draw[blue!35,dashed,rounded corners=3pt]
  (-1.65,-.15) rectangle (1.6,1.82);
 \draw[visible] (0,3.15)--(0,2.1);
 \draw[hidden] (0,2.1)--(-1.18,0);
 \draw[hidden] (0,2.1)--(1.18,0);
 \fill[black] (0,2.1) circle (1.7pt);
 \node[anchor=east,font=\footnotesize] at (-.13,2.32) {launch};
 \draw[visible] (2.2,3.15) .. controls (2.2,1.85) and (1.75,1.05)
  .. (.787,.7);
 \node[exposure point] at (.787,.7) {};
 \node[anchor=west,font=\footnotesize] at (1.5,.38) {first exposure};
 \node[align=center,font=\footnotesize] at (-1.18,-.48) {new label};
 \node[align=center,font=\footnotesize] at (1.18,-.48) {continued parent};
 \node[blue!65!black,font=\footnotesize] at (.05,-.98)
  {both descendant branches hidden};
\end{scope}
\begin{scope}[xshift=7cm]
 \node[font=\small\bfseries] at (.25,3.65) {Passage or virtual crossing};
 \fill[blue!4] (.18,-.15) rectangle (1.65,1.82);
 \draw[blue!35,dashed,rounded corners=3pt]
  (.18,-.15) rectangle (1.65,1.82);
 \draw[visible] (0,3.15)--(0,0);
 \draw[hidden] (0,2.1)--(1.18,0);
 \fill[black] (0,2.1) circle (1.7pt);
 \node[anchor=east,font=\footnotesize] at (-.13,2.32) {launch};
 \draw[visible] (2.2,3.15) .. controls (2.2,1.85) and (1.75,1.05)
  .. (.787,.7);
 \node[exposure point] at (.787,.7) {};
 \node[anchor=east,font=\footnotesize] at (-.06,-.48) {visible parent};
 \node[anchor=west,font=\footnotesize] at (.9,-.48) {new label};
 \node[blue!65!black,font=\footnotesize] at (.3,-.98)
  {only the new branch hidden};
\end{scope}
\end{tikzpicture}
\caption{The auxiliary exposure used to multiply contact savings.
Blue branches carry an unrevealed launch parameter. A genuine
scattering changes the continued parent's velocity, so both descendant
branches must be hidden. Passage leaves the parent unchanged.
Until first exposure, the black trajectories can be computed without
the hidden parameter. The construction is applied to exact signed
histories; no particles are deleted from the microscopic dynamics.
The drawing shows one selected tree. The full construction also omits
added groups whose positions have not yet been fixed relative to the
visible particles, including groups not selected for exposure;
Proposition~\ref{exp:exposure} gives the complete construction.}
\label{fig:exposure}
\end{figure}
